\maketitle

\begin{abstract}
本稿では、弱い重力波が生むジャイロスコピック重力メモリーのハード／ヘリシティ成分と、同じ局所放射域パッチの放射データを横切る固定した補助的ヌル Jacobi プローブ上の Bel--Robinson 曲率作用との間に鋭い下界を導く。対象は四次元、放射域の先頭次数、標準的な初期シアーなし Bondi フレーム、入出力定常条件を満たす線形化 TT/Bondi 放射データのクラスに限定する。ヌル Jacobi 散乱量 $\omega_\gamma$ と時間的世界線上のジャイロ向き角 $\Delta\Psi_{\rm H}$ は別個の観測量であるが、同じ局所放射データ上の二次汎関数として
\[
\omega_\gamma^{(2)}=-2\Delta\Psi_{\rm H}^{(2)},\qquad
|\omega_\gamma^{(2)}|=2|\Delta\Psi_{\rm H}^{(2)}|
\]
が成り立つ。

入出力の定常性は二偏光の潮汐制御に平均ゼロ条件を課す。このため物理的最小作用問題を
\[
\mathcal U_0=\left\{f\in L^2(0,1):\int_0^1 f(t)\,dt=0\right\},\qquad
K_{\rm stat}=P_0K_0P_0,
\]
\[
(K_0f)(t)=\frac12\int_0^1(t-s)|t-s|f(s)\,ds
\]
上の圧縮スペクトル問題として解く。ここで $P_0$ は $\mathcal U_0$ への直交射影である。定常 Jacobi 最小作用関数は
\[
\kappa_{\rm stat}(\omega)=\frac{|\omega|}{\|K_{\rm stat}\|}+O(\omega^2)
\]
を満たし、その係数は厳密終端回復によって鋭い。弱振幅族に対して
\[
Q_\gamma(\eps)=\eps^2Q_\gamma^{(2)}+O(\eps^3),\qquad
\Delta\Psi_{\rm H}(\eps)=\eps^2\Delta\Psi_{\rm H}^{(2)}+O(\eps^3)
\]
と二次摂動\emph{係数}を定義する。この leading-order 放射データ問題は純粋な二次変分問題であり、
\[
\boxed{\kappa_{\rm H}^{\rm lead}(\psi)=\frac{2}{\|K_{\rm stat}\|}|\psi|}
\]
と厳密に解ける。ここで \(\psi=\Delta\Psi_{\rm H}^{(2)}\) は full finite-amplitude memory ではなく、その二次摂動係数であり、\(\kappa_{\rm H}^{\rm lead}\) は \(Q_\gamma^{(2)}\) を最小化する。Gauss--Legendre Nystr\"om 法による検算では
\[
\|K_{\rm stat}\|=0.0281140680170\ldots,\qquad
\frac{2}{\|K_{\rm stat}\|}=71.13876223\ldots
\]
を得る。十進値は数値検証であり、定理の係数は作用素ノルムによって厳密に定義する。

ここで
\[
Q_\gamma=L^3\int_\gamma B(k,k,k,k)\,d\lambda
\]
は、固定した放射横断ヌル区間 $\gamma$ に付随する無次元の windowed Bel--Robinson 曲率作用である。これはジャイロの時間的世界線に沿う曝露量ではなく、また Bondi エネルギー、Isaacson エネルギー、放射エネルギー、検出器エネルギーでもない。固定区間のアフィン再パラメータ化には不変だが、区間端点そのものを変えれば別の最小化問題になる。また、異なる放射横断ヌル方向や異なるスクリーン同定を選べば別の補助プローブ問題であり、本稿は $Q_\gamma$ を方向独立な固有曝露量とはみなさない。本稿の主定理が最小化するのは、この固定ヌル窓上の量の弱振幅展開における先頭非零係数 $Q_\gamma^{(2)}$ である。本稿の主張はハード／ヘリシティ成分、弱振幅、放射域の先頭次数、初期シアーなしフレームに限定し、最終シアーは零に固定せず一定の displacement-memory 成分を許す。ジャイロスコピック・メモリー全体、スピン・メモリー、全 BMS フレームで不変な主張、有限振幅一様定理は主張しない。また、ここで構成する任意の回復波形が完全な漸近平坦真空 Einstein 時空へ大域的に埋め込めることも証明しない。
\end{abstract}

\section{はじめに}
重力波メモリーとは、一過性の重力放射が通過した後にも残る持続的な効果の総称である。変位メモリーに加え、スピン・メモリー、重心メモリー、ホロノミー、曲線偏差や回転する試験粒子に基づく持続的観測量など、複数の観測量が研究されている\cite{Flanagan2019,Flanagan2020,GrantNichols2022,SerajNeogi2023,Harte2025}。本稿が対象とするジャイロスコピック重力メモリーは、遠方の自由落下ジャイロが重力波通過後に獲得する向きの変化であり、先頭の $r^{-2}$ 項はソフト成分とハード成分に分かれる。既存研究ではソフト成分はスピン・メモリーと一致し、ハード成分は Bondi シアーに二次で、局所的な重力波ヘリシティおよび放射位相空間上の局所的な電気・磁気双対性の生成子として解釈される\cite{SerajOblak2023,FayeSeraj2025}。したがって、本稿はジャイロスコピック・メモリーやヘリシティを新たに発見するものではない。

本稿の問いは逆向きである。所定のハード・ジャイロスコピック・メモリーを持つ放射データに対し、その同じ波形を横切る固定ヌル窓上の Bel--Robinson 曲率作用はどこまで小さくできるかを問う。ここで資源として用いる $Q_\gamma$ は、ジャイロ世界線上の作用ではなく、補助的な放射横断ヌル測地線区間 $\gamma$ に沿う四重ヌル縮約の積分である。exact Jacobi 問題では $Q_\gamma$ 自体を用い、Bondi/TT 放射データとの比較では弱振幅展開の二次係数 $Q_\gamma^{(2)}$ を用いる。本稿で「鋭い」とは、単に一つの下界定数を与えるだけでなく、その定数より大きい普遍係数では不等式を維持できず、許容回復族によってその下界を任意の精度で回復できることを意味する。

本研究の数学的核は、二偏光のトレースレスなヌル Jacobi プロファイルに対するコンパクト反自己共役作用素
\[
(K_0f)(t)=\frac12\int_0^1(t-s)|t-s|f(s)\,ds
\]
である。無拘束の非線形 Jacobi メモリー・チャネルに対しては $\|K_0\|^{-1}$ が厳密な小目標の最良係数となる。しかし、物理的な入出力定常条件は潮汐制御に平均ゼロ条件を課すため、ハード・ジャイロスコピック・メモリー問題の最良係数は全 $L^2$ 空間のノルムではなく、平均ゼロ部分空間への圧縮
\[
K_{\rm stat}=P_0K_0P_0
\]
のノルムで決まる。本稿ではこの制約付き最小作用問題を解析的に解き、標準的な初期シアーなし Bondi フレームにおける Jacobi チャネルとハード・ジャイロスコピック・メモリーを結ぶ係数 2 の対応関係と組み合わせる。その結果、先頭次数のハード・メモリー二次係数を固定する限定放射データ問題に対して
\[
\kappa_{\rm H}^{\rm lead}(\psi)=\frac{2}{\|K_{\rm stat}\|}|\psi|
\]
という、$Q_\gamma^{(2)}$ に対する厳密で鋭い最小作用則を得る。exact nonlinear Jacobi endpoint の値関数は別に \(O(\omega^2)\) の補正を伴う小目標漸近則であり、両者を区別する。

既存研究との差は次のように限定する。Jacobi 伝播作用素、平面波メモリーテンソル、Bel--Robinson テンソル、ジャイロスコピック・メモリー、$N_{AB}\widetilde C^{AB}$ 型ハード汎関数、ヘリシティ／双対性の解釈は既知である\cite{Flanagan2020,Harte2025,SerajOblak2023,FayeSeraj2025}。本稿で新たに導く結果は、これらの既知構造に対して、同じ放射データ上のハード・ジャイロスコピック／ヘリシティ汎関数と固定ヌル窓上の Bel--Robinson 曲率作用を制約付き逆最適化として結び、その最良係数を確定する点にある。最終定理の同時条件に対する文献検索では対応する先行定理を特定できなかったが、この検索結果を歴史的優先権の証明とはみなさない。

入出力定常条件を課すと、無拘束問題の係数 $2/\|K_0\|=21.96874242\ldots$ は最良ではなくなる。定常条件から $\alpha,\beta\in\mathcal U_0$ が従い、この閉部分空間上で厳密終端展開とコンパクト作用素による回復構成をやり直すと、最良係数は $2/\|K_{\rm stat}\|$ となる。したがって旧係数は有効な弱い下界としては残るが、最良係数ではない。本稿の「鋭い」は、この制限作用素定理によって担保される。
